\section{Read the number of bytes the environment actually returned}
The file reader opens the terminated path with syscall 13 and read flag zero. A
negative result means the open failed. Otherwise retain the returned descriptor in a
saved register; it is a resource handle, not a byte pointer. The same descriptor is
passed to subsequent reads and the final close.

For syscall 14, place the descriptor in \texttt{a0}, the scratch address in
\texttt{a1}, and a positive request no larger than the scratch capacity in
\texttt{a2}. The result in \texttt{v0} determines the next action. A positive
result is the number of valid bytes placed at the start of the buffer. Zero means EOF.
A negative result is a read error. Test for the error before interpreting a result as
a length; a negative value is not a very large unsigned chunk.

Suppose the capacity is seven but a read returns three. Update with those three bytes.
The other four positions can still contain old workspace data. Hashing all seven
would append bytes that this read did not supply. After the update, issue another
read. A short positive result does not imply EOF, and a return of exactly the requested
capacity does not imply that another byte exists. EOF is established by a subsequent
zero result.

This project permits different positive request sizes within capacity. That freedom
allows equivalent buffering strategies while retaining the same file result. The
workspace itself is writable scratch storage; its contents are not a terminated string.
The low-level update routine still receives a precise read-only interval for each call.

The grader's explicit in-memory fixtures force short reads and failures at chosen
byte positions. This tests the control flow reproducibly without depending on disk
behavior or opening files on the instructor's machine. Ordinary IDE and CLI runs use
the existing MARS file environment. Their syscall result meanings are the same.
